\documentclass[12pt]{article}
%^ Start of document
\usepackage{times}  %Makes text TimesNewRoman
\usepackage{setspace} %Allows you to set single or double space 
\usepackage{biblatex} %Imports biblatex package
\usepackage{graphicx} % Required for inserting images
\usepackage{authblk}
\usepackage[colorlinks=true, urlcolor=blue, linkcolor=blue]{hyperref}
\usepackage{subfig}
\usepackage{float}
\usepackage{tabularx}
\usepackage{multirow}
\usepackage[paperheight=11in,
   paperwidth=8.5in,
   top=20mm,
   bottom=20mm,
   left=40mm,
   right=40mm]{geometry}
\usepackage{moresize}
\usepackage{tikz}
\usepackage{xcolor}
\usepackage{physics}
\usepackage{amssymb}
\usepackage{mathrsfs}
\usepackage{amsmath}
\usepackage[document]{ragged2e}
\usepackage{comment}
\usepackage{xcolor}
\addbibresource{Qhw.bib}
\begin{document}
\singlespacing  %Sets single spacing for whole document
\title{Planet Formation in Harsh Stellar Environments}

\author[1]{Zachary Tullier}
\affil[1]{Student\\Physics Department\\ University of Houston - Clear Lake \\Houston, Texas\\U.S}
\date{}
\maketitle

\section*{Abstract}
Dr. Megan Ryder’s lecture explores the process of star and planet formation, particularly within harsh radiation-rich environments like high-mass star-forming regions. She emphasizes how external influences---such as UV radiation from massive neighboring stars---can dramatically alter the chemistry, timescales, and survivability of planet-forming disks. The talk integrates observational data, theoretical modeling, and new insights from tools like the James Webb Space Telescope (JWST), focusing on how these factors shape planetary system architectures and the potential for habitability.

\section*{Introduction}
Planet formation is inherently linked to the environment in which stars are born. Dr. Ryder introduces the audience to the diversity of exoplanets discovered and compares these to our own solar system. Despite numerous observations, most known exoplanets do not resemble those of the solar system in mass or orbital separation. This mismatch reveals limitations in our detection technology and gaps in theoretical models.

\section*{Planet Formation: Ingredients and Environment}
Life as we understand it relies on liquid water and simple organic molecules---both influenced by a planet's location relative to its star and its chemical inventory. Ryder points out that the biochemistry of life on Earth depends on elements that are abundant in the universe, suggesting their presence may be common elsewhere. However, how planets acquire these ingredients and how they survive long enough to support life varies dramatically depending on their environment.

Key differences emerge when comparing:
\begin{itemize}
    \item \textbf{Terrestrial vs. Jovian planets}: Lower-mass, rocky planets are of greatest interest for life but harder to detect.
    \item \textbf{Planet-star distance}: A planet’s proximity to its host star affects its temperature, chemical composition, and formation potential.
    \item \textbf{Environmental radiation}: High-energy radiation from nearby massive stars may heat disks, destroy molecules, or alternatively stimulate new chemical pathways.
\end{itemize}

\section*{The Role of High-Mass Stars}
A central focus of the lecture is how high-mass stars, though few in number, dominate the radiation budget of star-forming regions and significantly influence their surroundings:
\begin{itemize}
    \item \textbf{Short lifespans} (a few million years) and intense UV radiation from these stars can erode planet-forming disks.
    \item \textbf{Star clusters like the Orion Nebula} and \textbf{Trumpler 14} provide direct observational evidence of disk destruction via photoevaporation.
    \item High UV flux shrinks disk size and mass, reducing the time available for planet formation and altering disk chemistry.
\end{itemize}

\section*{Observational Campaigns and Limitations}
Dr. Ryder underscores the difficulties of observing planet formation directly:
\begin{itemize}
    \item \textbf{Distance effects}: More distant regions, like the Carina Nebula ($\sim$7500 light years away), offer a rich variety of high-mass stars but are hard to resolve in detail.
    \item \textbf{Bias in target selection}: Telescopes often focus on large, bright disks (e.g., HL Tau) that may not represent typical planetary systems.
    \item \textbf{JWST contributions}: New observations are revealing both expected and unexpected chemical signatures. For example, CH$_3^+$ (a prebiotic molecule) is detected in high-radiation disks, while common biosignature precursors (like HCN or water) are often absent.
\end{itemize}

\section*{Case Study: The Tadpole Globule}
A highlight of the lecture is the analysis of a strange dusty structure dubbed the “Tadpole”:
\begin{itemize}
    \item It houses a young star protected within a dense globule, allowing cold gas to survive despite an intense external environment.
    \item Despite being ``boiled'' from the outside, the disk inside remains as cold as 15 K---ideal for some forms of planet chemistry.
    \item Spectroscopic evidence shows deuterated molecules, supporting the notion that protected regions can preserve complex chemistry even in extreme environments.
    \item This protected evolution stands in contrast to nearby stars exposed directly to radiation, which lose their disks quickly.
\end{itemize}

\section*{Ecosystem Effects and Theoretical Modeling}
Ryder proposes a shift from isolated disk models to \textbf{ecosystem thinking}, emphasizing:
\begin{itemize}
    \item \textbf{Cluster-wide simulations} show that UV radiation affects disk lifetimes non-linearly. A disk shielded for even 1 million years can form significantly more massive planets than one exposed earlier.
    \item \textbf{Age and shielding} interact complexly. Different clusters show vastly different disk survival rates depending on the number of nearby O-type stars.
    \item Ongoing studies (e.g., Sam Milstone’s work on M17) aim to map disk fractions as a function of age and radiation to constrain planetary formation timelines.
\end{itemize}

\section*{Future Directions}
Dr. Ryder outlines her team’s multi-pronged strategy:
\begin{enumerate}
    \item \textbf{Surveys} across multiple high-mass star-forming regions using integral field unit (IFU) spectroscopy to map stars, radiation fields, and surrounding nebulae.
    \item \textbf{Identification} of stars at various stages of exposure and their potential to host disks.
    \item \textbf{Selection} of targets for JWST and next-generation 30-meter telescopes to probe disk chemistry and evolution in real time.
\end{enumerate}

She stresses that understanding when and how stars are exposed to their environment is essential for understanding the diversity of planets.

\section*{Conclusion}
Dr. Ryder’s lecture emphasizes that planet formation is not a self-contained process, but one shaped by external radiation, disk chemistry, and cluster dynamics. Observations reveal that disks in high-radiation environments may follow unexpected paths---some losing water and organic precursors, while others reveal new chemical richness. Understanding these processes requires more observational data across a broader range of environments and timescales. Her research aims to fill this gap and better define the cosmic conditions that allow life to emerge.


\end{document}
